% !PRINTED 620
\weberpage{620}

\begin{center}
\S. 170.

\gesperrt{Minimalbasis des K\"orpers $\Omega_m$.}
\end{center}

Die S\"atze, die im vorigen Paragraphen bewiesen sind, f\"uhren uns zu dem folgenden Theorem:

II. \gesperrt{Die Zahlen}
\weq{1}{1, r, r^2, \ldots, r^{\mu-1}}
\gesperrt{bilden eine Basis des Systemes $\mathfrak{o}$ der ganzen Zahlen in $\Omega_m$.}

Um dies zu zeigen, gen\"ugt nach \S.\,145 der Nachweis, dass die ganze Zahl in $\Omega_m$
\weq{2}{\omega = x_0 + x_1 r + x_2 r^2 + \cdots + x_{\mu-1} r^{\mu-1},}
worin $x_0, x_1, \ldots, x_{\mu-1}$ ganze rationale Zahlen sind, nur dann durch eine rationale Primzahl $p$ theilbar sein kann, wenn die Zahlen $x_0, x_1, \ldots, x_{\mu-1}$ alle durch $p$ theilbar sind.

Nehmen wir also an, es sei $\omega$ durch $p$ theilbar; dann sind auch alle Zahlen $\omega_n$, die aus $\omega$ durch eine der $\mu$ Substitutionen $(r, r^n)$ entstehen, durch $p$ theilbar, und wir erhalten also aus (2) ein System von $\mu$ Gleichungen, die in Bezug auf die $x_i$ linear sind:
\weq{3}{\omega_n = x_0 + x_1 r^n + x_2 r^{2n} + \cdots + x_{\mu-1} r^{(\mu-1)n}.}

Die Determinante dieses Systemes, d.\,h. die aus den $\mu$ Reihen
\[ 1, r^n, r^{2n}, \ldots, r^{(\mu-1)n} \]
gebildete Determinante ist aber nach Bd. I, \S.\,46 gleich der Quadratwurzel $\sqrt{\varDelta}$, und wenn wir also das System der Gleichungen (3) aufl\"osen, so folgt, dass die s\"ammtlichen rationalen Zahlen $\varDelta x_i$ durch $p$ theilbar sein m\"ussen.

Ist nun $p$ nicht $=q$, so ist $\varDelta$ nicht durch $p$ theilbar, und s\"ammtliche $x_i$ m\"ussen durch $p$ theilbar sein.

Ist aber $p = q$, so setzen wir
\[ f(t) = x_0 + x_1 t + x_2 t^2 + \cdots + x_{\mu-1} t^{\mu-1}, \]
so dass $\omega_n = f(r^n)$ wird, und setzen darin $r = 1 - \sigma$ [\S.\,169, (6)]. Dann ist nach der \gesperrt{Taylor'schen} Entwickelung:
\weq{4}{\begin{aligned}
\omega &= f(1-\sigma) = f(1) - \sigma f'(1) + \dfrac{\sigma^2}{1.2} f''(1) - \cdots\\
&\qquad - \dfrac{\sigma^{\mu-1}}{\Pi(\mu-1)} f^{(\mu-1)}(1);
\end{aligned}}
%CONT>>
